\section{交织函子\index{jiaozhihanzi@交织函子|hyperpage}}
\label{sec:03.3.3}

设 $\theta$ 是 $\h^*$ 中的 Weyl 群轨道。若 $\theta$ 由正则元素构成，则由 1.12，范畴 $\D(U_\theta)$ 与范畴 $\D(\D_\lambda)$ 等价。特别地，对任意 $\lambda,\mu\in\theta$，范畴 $\D(\D_\lambda)$ 与 $\D(\D_\mu)$ 等价。这个等价由函子 $L\Delta_\mu\circ R\Gamma$ 给出，它从 $\D(\D_\lambda)$ 映入 $\D(\D_\mu)$。本节要用几何方式造一个与这个函子同构的函子。

先作一些几何准备。定义 $G$ 在 $X\times X$ 上的作用：
$$g\cdot(x,x')=(g\cdot x,\,g\cdot x'),\qquad g\in G,\quad (x,x')\in X\times X.$$
$X\times X$ 中的 $G$-轨道可按如下方式参数化。

先要在 $\g$ 的 Borel 子代数之间引入一种关系。设 $\bb$ 与 $\bb'$ 是 $\g$ 的两个 Borel 子代数，$\n$ 与 $\n'$ 是它们的幂零根，$N$ 与 $N'$ 是 $G$ 中对应的子群。设 $\mathfrak{c}$ 是 $\g$ 中含于 $\bb\cap\bb'$ 的 Cartan 子代数。记 $R$ 为 $(\g,\mathfrak{c})$ 在 $\mathfrak{c}^*$ 中的根系，$R^+$ 为 $\bb$ 决定的正根集。这给出 Cartan 三元组 $(\h^*,\Sigma,\Sigma^+)$ 到 $(\mathfrak{c}^*,R,R^+)$ 的一个特殊化。另一方面，$\bb'$ 在 $R$ 中决定另一组正根；经上述特殊化，这组正根对应某个唯一确定的 $w\in W$ 下的 $w(\Sigma^+)$。由于 $\mathfrak{c}$ 是 $\bb\cap\bb'$ 的 Levi 子代数，$\bb\cap\bb'$ 中一切 Cartan 子代数都被 $N\cap N'$ 的元素共轭。由此可见这个 $w\in W$ 不依赖 $\mathfrak{c}$ 的选取。于是我们说：$\bb'$ 相对于 $\bb$ 处于相对位置 $w$。设 $s:\h^*\to\mathfrak{c}^*$ 是 $\bb$ 决定的特殊化，则 $\bb'$ 决定的特殊化 $s'$ 满足 $s'=s\circ w$。由此可见 $\bb$ 相对于 $\bb'$ 处于相对位置 $w^{-1}$。

对 $w\in W$，置
$$Z_w=\{(x,x')\in X\times X\mid \bb_{x'}\ \text{相对于}\ \bb_x\ \text{处于相对位置}\ w\}.$$

\begin{lemma}\label{lem:03.3.1}
(i) 各 $Z_w$（$w\in W$）都是 $X\times X$ 的光滑子簇。
(ii) 映射 $w\mapsto Z_w$ 是 $W$ 到 $X\times X$ 中全体 $G$-轨道所成集合的双射。
\end{lemma}

\begin{proof}
固定 $w\in W$。集合 $Z_w$ 是 $G$-不变的，故含有一条 $G$-轨道 $O$。取 $x\in X$。因 $G$ 在 $X$ 上传递，$X\times X$ 中每条 $G$-轨道都与 $\{x\}\times X$ 相交，故存在 $x'\in X$ 使 $(x,x')\in O$。设 $(x,x'')\in Z_w$。取定 $\bb_x\cap\bb_{x'}$ 中的 Cartan 子代数 $\mathfrak{c}'$ 与 $\bb_x\cap\bb_{x''}$ 中的 Cartan 子代数 $\mathfrak{c}''$。则存在 $n\in N_x$ 使 $(\operatorname{Ad} n)(\mathfrak{c}')=\mathfrak{c}''$。因 $(x,x')$ 与 $(x,x'')$ 都在 $Z_w$ 中，故 $(\operatorname{Ad} n)(\bb_{x'})=\bb_{x''}$。于是 $x'$ 与 $x''$ 落在 $X$ 的同一 $B_x$-轨道内。这又推出 $O=Z_w$。
\end{proof}

记 $p_1$ 与 $p_2$ 分别为 $Z_w$ 到 $X\times X$ 两个因子的投影。

\begin{lemma}\label{lem:03.3.2}
纤维化 $p_i:Z_w\to X$（$i=1,2$）都局部平凡，纤维同构于 $\ell(w)$ 维仿射空间。投影 $p_i$（$i=1,2$）都是仿射态射。
\end{lemma}

\begin{proof}
只需讨论 $p_1$。设 $(x,x')\in Z_w$，并记 $B$、$B'$ 分别为 $x$、$x'$ 在 $G$ 中的稳定化子。则可用投影 $p_1:G/(B\cap B')\to G/B$ 代替 $p_1:Z_w\to X$。
\begin{annote}
底本此处引作「由 2」，与上下文不符：这里用到的是 3.1（$Z_w$ 是单条 $G$-轨道，故对 $(x,x')\in Z_w$ 有 $Z_w\simeq G/(B_x\cap B_{x'})$），而不是 3.2。已按 3.1 译。[来源：底本 p.48，3.1 的证明与轨道—稳定化子对应]
\end{annote}
设 $\overline{N}$ 是 $G$ 中与 $B$ 相对的 Borel 子群的幂幺根\index{miyaogen@幂幺根|hyperpage}。由 Bruhat 分解，自然映射 $\overline{N}\times B\to G$ 是到单位元某个开邻域的同构 [2, 14.13]。由此可见，轨道映射 $g\mapsto g\cdot x$ 给出 $\overline{N}$ 到 $x$ 在 $X$ 中某个开邻域 $U$ 的同构。此外，轨道映射 $g\mapsto g\cdot(x,x')$ 给出 $\overline{N}\times(B/(B\cap B'))$ 到 $p_1^{-1}(U)$ 的同构，使得图表
\begin{center}
\begin{tikzcd}
\overline{N}\times(B/(B\cap B')) \arrow[r] \arrow[d,"\mathrm{pr}_1"'] & p_1^{-1}(U) \arrow[d,"p_1"]\\
\overline{N} \arrow[r] & U
\end{tikzcd}
\end{center}
交换。$p_1$ 的纤维同构于 $B/(B\cap B')=N/(N\cap N')$，后者是维数为
$$\operatorname{Card}(\Sigma^+)-\operatorname{Card}(\Sigma^+\cap w(\Sigma^+))=\operatorname{Card}\bigl(\Sigma^+\cap(-w(\Sigma^+))\bigr)=\ell(w)$$
的仿射空间 [5, Ch. VI, §1, no. 6, Cor. 2 of Prop. 17]。
\end{proof}

设 $\Omega_{Z_w|X}$ 是投影 $p_1:Z_w\to X$ 的相对微分形式中最高次者构成的可逆 $\calO_{Z_w}$-模，$T_w$ 是它的逆模。$p_1$ 的纤维在 $(x,x')\in Z_w$ 处的切空间可等同于 $\n_x/(\n_x\cap\n_{x'})$，而 $\rho-w\rho$ 是 $\Sigma^+\cap(-w(\Sigma^+))$ 中诸根之和，故
$$T_w=p_1^*\bigl(\calO(\rho-w\rho)\bigr).$$

\begin{lemma}\label{lem:03.3.3}
(i) 设 $\nu\in P(\Sigma)$。则 $p_1^*(\calO(w\nu))=p_2^*(\calO(\nu))$。
(ii) 设 $\lambda\in\h^*$。则 $(\D_{w\lambda})_{p_1}=(\D^{p_2}_\lambda)_{T_w}$。
\end{lemma}

\begin{proof}
设 $(x,x')\in Z_w$。$(x,x')$ 在 $G$ 中的稳定化子 $B_x\cap B_{x'}$ 含 $G$ 的一个 Cartan 子群 $\ttt$，故 $\ttt$ 是它的 Levi 因子。由此可见 $B_x\cap B_{x'}$ 连通。显然有自然态射
$$\bb_x\cap\bb_{x'}\longrightarrow(\bb_x\cap\bb_{x'})/(\n_x\cap\bb_{x'})\longrightarrow\bb_x/\n_x$$
与
$$\bb_x\cap\bb_{x'}\longrightarrow(\bb_x\cap\bb_{x'})/(\bb_x\cap\n_{x'})\longrightarrow\bb_{x'}/\n_{x'},$$
而 $\bb_x/\n_x$ 与 $\bb_{x'}/\n_{x'}$ 都自然同构于 $\h$。这两个由 $\bb_x\cap\bb_{x'}$ 到 $\h$ 的自然态射相差 $w^{-1}$ 的作用。因此，齐性可逆 $\calO_{Z_w}$-模 $p_1^*(\calO(w\nu))$ 与 $p_2^*(\calO(\nu))$ 对应稳定化子的同一个特征标。这就证明了 (i)。

按定义，$(\D_{w\lambda})_{p_1}$ 与 $(\D^{p_2}_\lambda)_{T_w}$ 都是 $G$-齐性空间 $Z_w$ 上的 $G$-齐性扭微分算子层。如我们从……所知，$Z_w$ 上的 $G$-齐性扭微分算子层由 $\bb_x\cap\bb_{x'}$ 上 $B_x\cap B_{x'}$-不变的线性形式参数化。扭微分算子层 $(\D_{w\lambda})_{p_1}$ 对应 $w\lambda+\rho\in\h^*$ 经第一个态射在 $\bb_x\cap\bb_{x'}$ 上诱导的线性形式；扭微分算子层 $\D^{p_2}_\lambda$ 对应 $\lambda+\rho\in\h^*$ 经第二个同构诱导的线性形式。因此，要从 $\D^{p_2}_\lambda$ 得到 $(\D_{w\lambda})_{p_1}$，须用第一个同构下对应权
$$w\lambda+\rho-w(\lambda+\rho)=\rho-w\rho$$
的齐性可逆 $\calO_{Z_w}$-模去扭它。
\end{proof}

设 $w\in W$，$\lambda\in\h^*$。态射 $p_2:Z_w\to X$ 是满淹没，故 $p_2^+$ 是从 $\mathcal{M}_{qc}(\D_\lambda)$ 到 $\mathcal{M}_{qc}((\D_\lambda)_{p_2})$ 的正合函子（[8]，III.10.4.）。由 3.3，用 $T_w$ 去扭给出正合函子
$$V\longmapsto T_w\otimes_{\calO_{Z_w}}p_2^+(V),$$
它从 $\mathcal{M}_{qc}(\D_\lambda)$ 映入 $\mathcal{M}_{qc}((\D_{w\lambda})_{p_1})$。于是有函子 $V^\bullet\mapsto T_w\otimes_{\calO_{Z_w}}p_2^+(V^\bullet)$，它从 $\D^b(\D_\lambda)$ 映入 $\D^b((\D_{w\lambda})_{p_1})$。把它与直像\index{zhixiang@直像|hyperpage}函子 $Rp_{1+}:\D^b((\D_{w\lambda})_{p_1})\to\D^b(\D_{w\lambda})$ 复合，得到函子 $J_w:\D^b(\D_\lambda)\to\D^b(\D_{w\lambda})$，其公式为
$$J_w(V^\bullet)=Rp_{1+}\bigl(T_w\otimes_{\calO_{Z_w}}p_2^+(V^\bullet)\bigr),\qquad V^\bullet\in\D^b(\D_\lambda).$$

设 $V\in\mathcal{M}_{qc}(\D_\lambda)$。由 3.2，$p_1$ 是纤维维数为 $\ell(w)$ 的仿射态射，故当 $i<-\ell(w)$ 与 $i>0$ 时 $H^i(J_w(\D(V)))$ 为零（……）。
\begin{annote}
存疑：$p_1$ 的相对维数为 $\ell(w)$，故仿射态射的直像满足 $\mathrm{R}^ip_{1+}=0$ 对 $i>\ell(w)$ 成立（$i<0$ 时当然也为零）；底本此处作「$i<-\ell(w)$ 与 $i>0$」，上下界疑印反，疑本应作「$i<0$ 与 $i>\ell(w)$」。照底本译出，未改公式。[来源：底本 p.49；与同页 3.2「纤维为 $\ell(w)$ 维仿射空间」及仿射态射的消没定理对照]
\end{annote}
此外，函子
$$I_w(V)=R^0p_{1+}\bigl(T_w\otimes_{\calO_{Z_w}}p_2^+(V)\bigr)$$
从 $\mathcal{M}_{qc}(\D_\lambda)$ 映入 $\mathcal{M}_{qc}(\D_{w\lambda})$，并且右正合。我们称它为附着于 $w\in W$、在 $\mathcal{M}_{qc}(\D_\lambda)$ 与 $\mathcal{M}_{qc}(\D_{w\lambda})$ 之间的\emph{交织函子}。这样命名的理由后文自明。

\begin{lemma}\label{lem:03.3.4}
范畴 $\mathcal{M}_{qc}(\D_\lambda)$ 有足够多的投射对象。
\end{lemma}

\begin{proof}
用扭操作显然可设 $\lambda$ 反支配且正则。在这种情形，由 1.3，$\mathcal{M}_{qc}(\D_\lambda)$ 等价于 $\mathcal{M}(U_\theta)$。
\end{proof}

因此，对任意 $w\in W$，可定义 $I_w$ 的左导出函子
$$LI_w:\D^-(\D_\lambda)\longrightarrow\D^-(\D_{w\lambda}).$$
后文将看到，这个函子限制在 $\D^b(\D_\lambda)$ 上与 $J_w$ 一致。

下面研究这些函子的若干基本性质。先分析它们在几何平移下的行为。

\begin{lemma}\label{lem:03.3.5}
设 $w\in W$，$\lambda\in\h^*$，$\nu\in P(\Sigma)$。则
\begin{enumerate}[label=(\roman*),leftmargin=2.2em,itemsep=0pt,topsep=2pt]
\item 对任意 $V^\bullet\in\D^b(\D_\lambda)$ 有 $J_w(V^\bullet(\nu))=J_w(V^\bullet)(w\nu)$；
\item 对任意 $V^\bullet\in\D^-(\D_\lambda)$ 有 $LI_w(V^\bullet(\nu))=LI_w(V^\bullet)(w\nu)$。
\end{enumerate}
\end{lemma}

\begin{proof}
先证第一个关系。由 3.3(i)，对任意 $V\in\mathcal{M}_{qc}(\D_\lambda)$ 有
\begin{align*}
p_2^+\bigl(V(\nu)\bigr)&=p_2^+\bigl(V\otimes_{\calO_X}\calO(\nu)\bigr)
=p_2^+(V)\otimes_{\calO_{Z_w}}p_2^*(\calO(\nu))\\
&=p_2^+(V)\otimes_{\calO_{Z_w}}p_1^*(\calO(w\nu))
\end{align*}
作为 $(\D_{\lambda+\nu})_{p_2}$-模。因直像函子关于目标簇是局部的，
\begin{align*}
J_w(V^\bullet(\nu))&=Rp_{1+}\bigl(T_w\otimes_{\calO_{Z_w}}p_2^+(V^\bullet(\nu))\bigr)\\
&=Rp_{1+}\bigl(T_w\otimes_{\calO_{Z_w}}p_2^+(V^\bullet)\otimes_{\calO_{Z_w}}p_1^*(\calO(w\nu))\bigr)\\
&=Rp_{1+}\bigl(T_w\otimes_{\calO_{Z_w}}p_2^+(V^\bullet)\bigr)\otimes_{\calO_X}\calO(w\nu)\\
&=J_w(V^\bullet)(w\nu)
\end{align*}
对任意 $V^\bullet\in\D^b(\D_\lambda)$ 成立。特别地，对任意 $V\in\mathcal{M}_{qc}(\D_\lambda)$ 有 $I_w(V(\nu))=I_w(V)(w\nu)$。由于扭保持投射对象，引理得证。
\end{proof}

下一步是函子 $J_w$（$w\in W$）的「乘积公式」。先要对 $X\times X$ 中的 $G$-轨道补充一些几何信息。设 $w,w'\in W$。记 $p_1,p_2$ 为 $Z_w$ 到 $X$ 的投影，$p'_1,p'_2$ 为 $Z_{w'}$ 到 $X$ 的相应投影。设 $Z_{w'}\times_X Z_w$ 是 $Z_{w'}$ 与 $Z_w$ 关于态射 $p'_2$ 与 $p_1$ 的纤维积。记 $q':Z_{w'}\times_X Z_w\to Z_{w'}$ 与 $q:Z_{w'}\times_X Z_w\to Z_w$ 为到第一、第二因子的相应投影。由 3.2，$q$ 与 $q'$ 都是仿射态射。最后，态射 $p'_1\circ q':Z_{w'}\times_X Z_w\to X$ 与 $p_2\circ q:Z_{w'}\times_X Z_w\to X$ 确定一个态射 $r:Z_{w'}\times_X Z_w\to X\times X$。于是有下述交换图表。
\begin{center}
\begin{tikzcd}[column sep=3.6em,row sep=2.9em]
Z_{w'}\times_X Z_w \arrow[r,"q"] \arrow[d,"q'"'] \arrow[rr,bend left=18,"r"] & Z_w \arrow[dd,"p_1"'{pos=0.78}] \arrow[dr,"p_2"] & X\times X\\
Z_{w'} \arrow[rr,"p'_1"'{pos=0.2}] \arrow[dr,"p'_2"'] & & X\\
{} & X & {}
\end{tikzcd}
\end{center}
此外，图表中一切态射都是 $G$-等变的。由构造可见 $r$ 的像含于 $Z_{w'w}$。故由 $r$ 的 $G$-等变性，$r$ 是 $Z_{w'}\times_X Z_w$ 到 $Z_{w'w}$ 的满射。

\begin{lemma}\label{lem:03.3.6}
设 $w,w'\in W$ 满足 $\ell(w'w)=\ell(w')+\ell(w)$。则 $r:Z_{w'}\times_X Z_w\to Z_{w'w}$ 是同构。
\end{lemma}

\begin{proof}
由 3.2 知
$$\dim\,(Z_{w'}\times_X Z_w)=\dim X+\ell(w)+\ell(w')=\dim X+\ell(w'w),$$
而 $\dim Z_{w'w}=\dim X+\ell(w'w)$。由 $r$ 的 $G$-等变性，$Z_{w'}\times_X Z_w$ 中每条 $G$-轨道 $O$ 都满射到 $Z_{w'w}$。故
$$\dim O=\dim X+\ell(w'w)=\dim\,(Z_{w'}\times_X Z_w),$$
即 $O$ 在 $Z_{w'}\times_X Z_w$ 中开。另一方面，$Z_{w'}\times_X Z_w$ 不可约，故 $O=Z_{w'}\times_X Z_w$。这说明 $Z_{w'}\times_X Z_w$ 是覆叠 $Z_{w'w}$ 的 $G$-齐性空间。由于 $Z_{w'w}$ 中基点的稳定化子连通，$r$ 是同构。
\end{proof}

因此，若 $w,w',w''\in W$ 满足 $w''=w'w$ 且 $\ell(w'')=\ell(w')+\ell(w)$，就可以把 $Z_{w''}$ 与 $Z_{w'}\times_X Z_w$ 等同。在这个等同下，$Z_{w''}$ 到 $X$ 的投影 $p''_1$、$p''_2$ 分别对应映射 $p'_1\circ q'$ 与 $p_2\circ q$。此外还有下述结果。

\begin{lemma}\label{lem:03.3.7}
设 $w,w'\in W$ 满足 $\ell(w'w)=\ell(w')+\ell(w)$。则 $J_{w'}\circ J_w=J_{w'w}$。
\end{lemma}

\begin{proof}
令 $w''=w'w$。由 3.3(i) 有
{\small
\begin{align*}
&q'^*(T_{w'})\otimes_{\calO_{Z_{w'}\times_X Z_w}}q^*(T_w)\\
&\ = (p'_1\circ q')^*\bigl(\calO(\rho-w'\rho)\bigr)\otimes_{\calO_{Z_{w'}\times_X Z_w}}(p_1\circ q)^*\bigl(\calO(\rho-w\rho)\bigr)\\
&\ = (p'_1\circ q')^*\bigl(\calO(\rho-w'\rho)\bigr)\otimes_{\calO_{Z_{w'}\times_X Z_w}}(p'_2\circ q')^*\bigl(\calO(\rho-w\rho)\bigr)\\
&\ = q'^*\Bigl(p'^*_1\bigl(\calO(\rho-w'\rho)\bigr)\otimes_{\calO_{Z_{w'}}}p'^*_2\bigl(\calO(\rho-w\rho)\bigr)\Bigr)\\
&\ = q'^*\Bigl(p'^*_1\bigl(\calO(\rho-w'\rho)\bigr)\otimes_{\calO_{Z_{w'}}}p'^*_1\bigl(\calO(w'\rho-w'w\rho)\bigr)\Bigr)\\
&\ = q'^*\bigl(p'^*_1(\calO(\rho-w''\rho))\bigr)=(p'_1\circ q')^*\bigl(\calO(\rho-w''\rho)\bigr)=T_{w''},
\end{align*}}
这里的等号是在 $Z_{w'}\times_X Z_w$ 与 $Z_{w''}$ 的等同下取的。于是由基变换\index{jibianhuan@基变换|hyperpage}（……），
{\small
\begin{align*}
J_{w'}(J_w(V^\bullet))
&=Rp'_{1+}\Bigl(T_{w'}\otimes_{\calO_{Z_{w'}}}p'^+_2\bigl(Rp_{1+}(T_w\otimes_{\calO_{Z_w}}p_2^+(V^\bullet))\bigr)\Bigr)\\
&=Rp'_{1+}\Bigl(T_{w'}\otimes_{\calO_{Z_{w'}}}Rq'_{+}\bigl(q^+(T_w\otimes_{\calO_{Z_w}}p_2^+(V^\bullet))\bigr)\Bigr)\\
&=R(p'_1\circ q')_{+}\Bigl(q'^*(T_{w'})\otimes_{\calO_{Z_{w'}\times_X Z_w}}(p_2\circ q)^+(V^\bullet)\Bigr)\\
&=J_{w''}(V^\bullet)
\end{align*}}
对任意 $V^\bullet\in\D^b(\D_\lambda)$ 成立。
\end{proof}

\begin{corollary}\label{cor:03.3.8}
设 $w,w'\in W$ 满足 $\ell(w'w)=\ell(w')+\ell(w)$。则 $I_{w'}\circ I_w=I_{w'w}$。
\end{corollary}

下面要更细致地分析附着于单反射的函子。固定单根 $\alpha\in\Pi$。为简化后文的记号，置 $Z=Z_{s_\alpha}$，$T=T_{s_\alpha}$。此时由 3.2，投影 $p_1:Z\to X$ 的纤维是仿射线。于是可把 $T$ 看作 $\calO_Z$ 上由切于 $p_1$ 纤维的局部向量场构成的可逆模\index{kenimo@可逆模|hyperpage}。

\begin{lemma}\label{lem:03.3.9}
设 $\alpha\in\Pi$，$\lambda\in\h^*$。则当 $i\neq 0$ 时 $H^i(J_{s_\alpha}(\D(\D_\lambda)))=0$。
\end{lemma}

3.9 的证明将是下述讨论的推论；这段讨论还会给出关于交织函子 $I_{s_\alpha}$ 的作用的更细致信息。

基本想法是把分析化归到 $\g=\mathfrak{sl}(2,\C)$ 的情形。这一化归基于一个局部平凡化结果。设 $X_\alpha$ 是全体 $\alpha$ 型抛物子代数\index{paowuzidaishu@抛物子代数|hyperpage}构成的广义旗簇。$\g$ 中任一 Borel 子代数 $\bb_x$ 都含于唯一的 $\alpha$ 型抛物子代数 $\p_y$（$y\in X_\alpha$），故有自然投影 $p_\alpha:X\to X_\alpha$。对任意 $x\in X$，纤维 $p_\alpha^{-1}(p_\alpha(x))$ 由 $x$ 以及一切满足「$\bb_{x'}$ 相对于 $\bb_x$ 处于相对位置 $s_\alpha$」的 $x'\in X$ 构成。

\begin{lemma}\label{lem:03.3.10}
投影 $p_\alpha:X\to X_\alpha$ 局部平凡，其纤维同构于射影直线 $\mathbb{P}^1$。
\end{lemma}

\begin{proof}
固定 $y\in X_\alpha$ 与 $x\in X$，使 $p_\alpha(x)=y$。记 $B$、$P$ 分别为 $x$、$y$ 在 $G$ 中的稳定化子。设 $P'$ 是与 $P$ 相对的抛物子群\index{paowuziqun@抛物子群|hyperpage}，$N'$ 是它的幂幺根。则自然映射 $N'\times P\to G$ 是到 $G$ 中单位元某个开邻域的同构（[?]，4.2）。因此，由轨道映射 $g\mapsto g\cdot y$ 诱导的 $N'\to X_\alpha$ 的自然态射是 $N'$ 到 $y$ 的某个开邻域 $U$ 的同构。此外，轨道映射 $g\mapsto g\cdot x$ 给出 $N'\times(P/B)$ 到 $p_\alpha^{-1}(U)$ 的同构，使得图表
\begin{center}
\begin{tikzcd}
N'\times(P/B) \arrow[r] \arrow[d,"\mathrm{pr}_1"'] & p_\alpha^{-1}(U) \arrow[d,"p_\alpha"]\\
N' \arrow[r] & U
\end{tikzcd}
\end{center}
交换。这说明纤维同构于 $P/B$。设 $R$ 是 $P$ 的根。则 $R\subset B$，故 $P/B=(P/R)/(B/R)$。由于 $P/R$ 是 $PSL(2,\C)$ 的覆叠，而 $B/R$ 是它的 Borel 子群，$P/B$ 同构于 $\mathbb{P}^1$。
\end{proof}
\begin{annote}
底本此处文献编号缺字，印作「([?], 4.2)」，其中「?」为底本缺失的编号。因无法确知原编号，照录为「[?]」，未擅自补。[来源：底本 p.52；同一事实在 3.2 的证明中引作 [2, 14.13]]
\end{annote}

我们指出：$p_\alpha^{-1}(U)$ 是 $P'$ 的齐性空间；若记 $L$ 为 $P'$ 与 $P$ 的公共 Levi 因子，则 $p_\alpha^{-1}(U)$ 可与 $N'\times(L/(L\cap B))$ 等同。设 $M$ 是 $L$ 关于其中心的商，$\m$ 是它的 Lie 代数，$X_\m$ 是相应的旗簇。显然 $M$ 同构于 $PSL(2,\C)$，$X_\m$ 同构于 $\mathbb{P}^1$。在 $X$ 与 $X_\m$ 中分别取基点 $\bb$ 与 $\bb\cap\m$，就确定 Cartan 代数 $\h_\m$ 到 $\h$ 的一个自然包含；它把 $\h^*_\m$ 中的根系 $\Sigma_\m$ 与 $\alpha$、$-\alpha$ 的限制等同起来，而 $\Sigma_\m$ 中的正根 $\beta$ 对应 $\alpha$。借助映射 $\mu\mapsto\check{\beta}(\mu)$，可把 $\h_\m$ 的对偶空间与 $\C$ 等同。由齐性扭微分算子层的讨论（……）可见，对任意 $\lambda\in\h^*$ 有
$$\D_\lambda|_{p_\alpha^{-1}(U)}=\D_{N'}\boxtimes\D_{\check{\alpha}(\lambda)},$$
这里 $\D_{\check{\alpha}(\lambda)}$ 记 $\mathbb{P}^1=X_\m$ 上由 $\check{\alpha}(\lambda)\in\C$ 按上述对应确定的齐性扭层。按定义 $p_\alpha\circ p_1=p_\alpha\circ p_2$，故
$$p_1^{-1}\bigl(p_\alpha^{-1}(U)\bigr)=p_2^{-1}\bigl(p_\alpha^{-1}(U)\bigr),$$
作为 $Z$ 的开子集。此外，在上述等同下，它与 $N'\times((X_\m\times X_\m)-\Delta_\m)$ 同构，其中 $\Delta_\m$ 记 $X_\m\times X_\m$ 中的对角线。设 $q_1$、$q_2$ 分别是簇 $Z_\m=(X_\m\times X_\m)-\Delta_\m$ 到 $X_\m$ 的第一、第二因子投影诱导的态射。则利用上述等同，
\begin{align*}
&R^ip_{1+}\bigl(T\otimes_{\calO_Z}p_2^+(\D_\lambda)\bigr)|_{p_\alpha^{-1}(U)}\\
&\qquad=\D_{N'}\boxtimes R^iq_{1+}\bigl(T_\m\otimes_{\calO_{Z_\m}}q_2^+(\D_{\check{\alpha}(\lambda)})\bigr),
\end{align*}
其中 $T_\m$ 记 $Z_\m$ 上由切于 $q_1$ 纤维的局部向量场构成的可逆 $\calO_{Z_\m}$-模。

现在可以证明 3.9。上述讨论把计算化归到 $\g=\mathfrak{sl}(2,\C)$ 的情形。以下就设 $\g=\mathfrak{sl}(2,\C)$。显然 $\D_\lambda$ 是 $G$-齐性 $\calO_X$-模，故 $T\otimes_{\calO_Z}p_2^+(\D_\lambda)$ 是 $G$-齐性 $\calO_Z$-模。它在 $p_1$ 下的直像是 $G$-齐性 $\calO_X$-模。因此，作为 $\calO_X$-模，它们完全由基点 $x\in X$ 处的几何纤维决定；而它们的高次几何纤维为零。令 $F=X-\{x\}$。则 $p_1^{-1}(x)$ 是 $Z$ 的光滑闭子簇，等于 $\{x\}\times F$。设 $i_F:\{x\}\times F\to Z$ 与 $i_x:\{x\}\to X$ 是自然包含，$r$ 是 $\{x\}\times F$ 到 $x$ 的投影。则有下述交换图表：
\begin{center}
\begin{tikzcd}
\{x\}\times F \arrow[r,"i_F"] \arrow[d,"r"'] & Z \arrow[d,"p_1"]\\
\{x\} \arrow[r,"i_x"'] & X
\end{tikzcd}
\end{center}
